%% Macros de dados do documento
\autor{Autor do projeto}
\titulo{Título do projeto}
\instituicao{Instituto Federal do Espírito Santo}
\curso{Bacharelado em Sistemas de Informação}
\local{Cachoeiro de itapemirim}
\data{2026}


\tipotrabalho{Trabalho de Graduação}
\preambulo{Trabalho de Conclusão de Curso apresentado à Coordenadoria do Curso de Sistemas de Informação do
Instituto Federal de Educação, Ciência e Tecnologia do Espírito Santo, como requisito parcial para a obtenção do título de Bacharel em Sistemas de Informação.}

\orientador[Orientador][Prof. Dr.]{Beltrano}[de Tal]{Instituto Federal do Espírito Santo}
\coorientador[Coorientadora][Profa. Dra.]{Beltrana}[de Tal]{Instituto Federal do Espírito Santo}

\examinadori{Profa. Dra. Fulana de Tal}{Instituto Federal do Espírito Santo}{Examinadora}
\examinadorii{Prof. Dr. Cicrano de Tal}{Instituto Federal do Espírito Santo}{Examinador}

\approvaldate{01}{Agosto}{2026}

% Palavras-chaves
\palavraschave{Palavra Chave 1}[Palavra Chave 2][Palavra Chave 3][Palavra Chave 4][Palavra Chave 5]
\keywords{keyword 1}[keyword 2][keyword 3][keyword 4][keyword 5]

% Ficha catalográfica
\fichabiblioteca{
    Dados Internacionais de Catalogação-na-Publicação (CIP) \\
    (Biblioteca Nilo Peçanha do Instituto Federal do Espírito Santo)
}
\fichaautor{Elaborada por XXXXXXXXXXXXXXXXXXXXX – CRB-X/ES - XXX}
\fichatags{
    1. Sistemas de Informação.  
    2. Redes de Computadores – Protocolos.  
    3. Segurança da Informação – Criptografia.  
    4. Computação em Nuvem – Serviços.  
    5. Internet das Coisas (IoT).  
    6. Engenharia de Software – Desenvolvimento de Sistemas.  
    I. Borges, Everson Scherrer.  
    II. Instituto Federal do Espírito Santo.  
    III. Título.
}

\cutter{X999y}
\cdd{000.00}
\selectlanguage{brazil}

% Arial (para usar times-new-roman, comente as três linhas abaixo)
\usepackage{helvet}
\renewcommand{\familydefault}{\sfdefault}
\renewcommand{\ABNTEXchapterfont}{\bfseries}
% Comandos para revisar o texto
\usepackage{easyReview} 


%% Elementos opcionais
\editardedicatoria{\input{pre_textuais/dedicatoria}}
\editaragradecimentos{\input{pre_textuais/agradecimentos}}
\editarepigrafe{\input{pre_textuais/epigrafe}}


%% Listas [Siglas e Simbolos]
\editarlistasiglas{\input{pre_textuais/siglas}}
\editarlistasimbolos{\input{pre_textuais/simbolos}}